% concatenated from CullisGeneralPart2toArXiV.tex
\documentclass{elsarticle}

\usepackage[utf8]{inputenc}
%\usepackage{enumerate}
\usepackage{amsmath,amssymb,amsthm,amsfonts}
\usepackage{bm}
\usepackage[hidelinks,colorlinks = true]{hyperref}
\usepackage{enumitem}
\usepackage{blkarray}
\setenumerate{labelindent=0.25\parindent,leftmargin=*}
%\usepackage[a4paper, total={16.5cm, 24cm}]{geometry}
\setcounter{MaxMatrixCols}{12}
\usepackage{enumitem}
\usepackage{blkarray}
\usepackage{tikz-cd}
\usepackage{tcolorbox}
\newtcolorbox{mybox}{colback=red!5!white,colframe=red!75!black, sharp corners = all}
\usepackage{bclogo} %\bcpanchant
\def\ucsign{\bcpanchant}
%\newsavebox{\mybox}
\newenvironment{underconst}{\begin{mybox}\par\noindent\hangindent=2em\hangafter=-2
  \smash{\hbox to0pt{\hskip-\hangindent
    \lower1.5\baselineskip\hbox{\ucsign}\hfill}}
    }{\end{mybox}}
    
\newcommand{\tunderbrace}[1]{\underbrace{\textstyle#1}}
    
%\newenvironment{underconst}{\noindent\begin{lrbox}{\mybox}\begin{minipage}{\textwidth}\par\noindent\hangindent=2em\hangafter=-2
%  \smash{\hbox to0pt{\hskip-\hangindent
%    \lower1.5\baselineskip\hbox{\ucsign}\hfill}}
%    }{\end{minipage}\end{lrbox}\fbox{\usebox{\mybox}}}

  %  {\lower\baselineskip\hbox{\ucsign}}
\DeclareMathOperator{\rk}{rk}
\DeclareMathOperator{\id}{id}
\DeclareMathOperator{\cl}{cl}
\DeclareMathOperator{\add}{add}
\DeclareMathOperator{\rad}{rad}
\DeclareMathOperator{\ins}{ins}
\DeclareMathOperator{\del}{del}
\DeclareMathOperator{\Ker}{ker}
%\DeclareMathOperator{\ker}{ker}
\DeclareMathOperator{\Img}{Im}
%\DeclareMathOperator{\nul}{Im}
\DeclareMathOperator{\Coimg}{Coim}
\DeclareMathOperator{\codim}{codim}
\DeclareMathOperator{\Char}{char}
\DeclareMathOperator{\Span}{span}
\DeclareMathOperator{\diag}{diag}
\DeclareMathOperator{\colsp}{\mathsf C}
\DeclareMathOperator{\rowsp}{\mathsf R}
\DeclareMathOperator{\nullsp}{\mathsf N}

\newcommand{\m}{\mathbf M}
\def\F{{\mathbb F}}
\def\K{{\mathsf K}}
%\def\L{{\mathsf L}}
\def\M{{\mathcal M}}
\def\N{{\mathbb N}}
\def\SCS{{\mathbf S}}
%\DeclareMathOperator{\ZC}{C^{+}}
%\DeclareMathOperator{\DC}{C^{--}}



%\DeclareMathOperator{\Det}{Det}
\let\det\undefined
\DeclareMathOperator{\det}{det}
\DeclareMathOperator{\Det}{Det}
\DeclareMathOperator{\dc}{d_C}
\DeclareMathOperator{\tr}{tr}
\DeclareMathOperator{\sgn}{sgn}
\DeclareMathOperator{\row}{row}
\DeclareMathOperator{\col}{col}
\DeclareMathOperator{\Hom}{Hom}
\DeclareMathOperator{\ALS}{\mathbf{ALS}}
\DeclareMathOperator{\AS}{\mathcal{F}}

\newenvironment{psmallmatrix}
  {\left(\begin{smallmatrix}}
  {\end{smallmatrix}\right)}






\theoremstyle{plain}
\newtheorem{theorem}{Theorem}[section]
\newtheorem*{theorem*}{Theorem}

\newtheorem{lemma}[theorem]{Lemma}
\newtheorem{claim}[theorem]{Claim}
%\newtheorem{cor}[theorem]{Corollary}
%\newtheorem{cor}[thm]{Corollary}
%\newtheorem{claim}[theorem]{claim}
\newtheorem{proposition}[theorem]{Proposition}
\newtheorem{corollary}[theorem]{Corollary}
\newtheorem{conjecture}[theorem]{Conjecture}


%\newtheorem{prop}[thm]{Proposition}


\theoremstyle{definition}
\newtheorem{definition}[theorem]{Definition}
\newtheorem*{definition*}{Definition}
\newtheorem{question}[theorem]{Question}
%\newtheorem{defi}[thm]{Definition}
\newtheorem{example}[theorem]{Example}
\newtheorem*{example*}{Example}
\newtheorem{counterexample}[theorem]{Counterexample}
% \newtheorem{exm}[thm]{Example}
\newtheorem{problem}[theorem]{Problem}
%\newtheorem{problem}[thm]{Problem}
\newtheorem{remark}[theorem]{Remark}
\newtheorem{note}[theorem]{Note}
%\newtheorem{rem}[thm]{Remark}
\numberwithin{equation}{section}


\begin{document}




\title{Linear maps preserving the Cullis'
determinant. II}
\author{Alexander Guterman}
\ead{alexander.guterman@biu.ac.il}
\author{Andrey Yurkov\corref{cor1}}
\ead{andrey.yurkov@biu.ac.il}
\cortext[cor1]{Corresponding author}

\affiliation{organization={Department of Mathematics, Bar Ilan University},addressline={Ramat-Gan},postcode={5290002}, country={Israel}}

\begin{keyword}
\MSC[2020] 15A15 \sep 15A86 \sep 47B49\\
Cullis' determinant \sep linear preservers \sep rectangular matrices    
\end{keyword}

\begin{abstract}
This paper is the second in the series of papers devoted to the explicit description of linear maps preserving the Cullis' determinant of rectangular matrices with entries belonging to an arbitrary ground field which is large enough.

In this part we solve the linear preserver problem for the Cullis' determinant defined on the spaces of matrices of size $n\times k$ with $k \ge 4,\; n \ge k + 2$ and $n + k$ is odd. 

In comparison with the case when $n + k$ is even, in this case linear maps preserving the Cullis' determinant could be singular and are represented as a sum of two linear maps: first is two-sided matrix multiplication and second is any linear map whose image consists of matrices, all rows of which are equal.
\end{abstract}



\maketitle



\section{Introduction}
The theory of linear maps preserving matrix invariants is well-known direction of research in Linear Algebra connected with other fields of mathematics. It has a long history which dates back to the beginning of 20th century in the works of Frobenius~\cite{GF} and continues at present time (see~\cite{LAMA199233} for the survey of the results until the end of 20th century and~\cite{hogben_handbook_2014} for a brief introduction).

The Cullis' determinant is an invariant of rectangular matrix generalizing the notion of the ordinary determinant of square matrix which was introduced by Cullis in~\cite{cullis1913} and studied later by Radi\'{c} in~\cite{radic1966,radic2005,radic1991,radic2008}, Makarewicz, Mozgawa, Pikuta and Sza\l{}kowski in~\cite{makarewicz2014,makarewicz2016,makarewicz2020}, Amiri, Fathy and Bayat in~\cite{amiri2010}, and Nakagami and Yanay in~\cite{NAKAGAMI2007422}. Thus, the question regarding the description of linear maps preserving the Cullis' determinant arises naturally. The detailed introduction to the subject could be found in~\cite{Guterman2025}. %, which is the first paper in these series of papers.

In particular, \cite{Guterman2025} contains the description of linear maps preserving the Cullis determinant $\det_{n\,k}$ for the case when $ k \ge 4$, $n \ge k + 2$ and $n + k$ is even as it stated in the following theorem (all necessary definitions are given in Section~\ref{sec:prelim}).

\begin{theorem}[{\cite[Theorem~5.14]{Guterman2025}}]\label{MainTheoremEvenIntro}Assume that $k \ge 4, n \ge k + 2, n + k$ is even and $|\F| > k$. Let $T\colon \M_{n\,k} (\F) \to \M_{n\,k} (\F)$ be a linear map. Then $\det_{n\,k} (T(X)) = \det_{n\,k}(X)$ for all $X \in \M_{n\,k} (\F)$ if and only if there exist $A \in \M_{n\,n}(\F)$ and $B \in \M_{k\,k}(\F)$ such that
\[
\det_{n\,k} \Bigl(A(|i_1,\ldots, i_k]\Bigr) \det_k \Bigl(B\Bigr) = (-1)^{i_1 + \ldots + i_k - 1 - \ldots - k}
\]
for all increasing sequences $1 \le i_1 < \ldots < i_k \le n $ and
\[
T(X) = AXB
\]
for all $X \in \M_{n\,k} (\F).$
\end{theorem}

The parity of $n + k$ is significant because if $n + k$ is odd, then there exist linear maps preserving the Cullis' determinant which do not have such form (Corollary~\ref{cor:AdditionLinearMapNPLusKOdd}). 

However, in this paper we show that the case $(n,k)$ with $n \ge k$ and $n + k$ odd reduces to the case $(n-1,k)$ (Lemma~\ref{lem:FromNKOddToNKEvenIfKerW}\ref{lem:FromNKOddToNKEvenIfKerW:part6}). Using this reduction we prove the main theorem of this paper providing the description of linear maps preserving $\det_{n\,k}$ if $n + k$ is odd.


\begin{theorem}[Theorem~\ref{thm:MainTheoremCullisNKOdd}]Assume that $|\F| > k \ge 4$, $n \ge k + 2$ and $n + k$ is odd. Let $T\colon \M_{n\,k} (\F) \to \M_{n\,k} (\F)$ be a linear map. Then $\det_{n\, k} (T(X)) = \det_{n\,k}(X)$ for all $X \in \M_{n\,k} (\F)$ if and only if  there exist $A \in \M_{n\, n}(\F)$ and $B \in \M_{k\, k}(\F)$ such that
\[
\det_{n\,k} \Bigl(A(|i_1,\ldots, i_k]\Bigr) \det_k \Bigl(B\Bigr) = (-1)^{i_1 + \ldots + i_k - 1 - \ldots - k}
\]
for all increasing sequences $1 \le i_1 < \ldots < i_k \le n $ and a linear map\linebreak $\phi\colon \M_{n\,k}(\F) \to W_{n\,k}$ such that
\[
\det_{n\,k} \Bigl(A(|i_1,\ldots, i_k]\Bigr) \det_k \Bigl(B\Bigr) = (-1)^{i_1 + \ldots + i_k - 1 - \ldots - k}
\]
for all increasing sequences $1 \le i_1 < \ldots < i_k \le n $ and
\begin{equation*}
T(X) = AXB + \phi(X)
\end{equation*}
for all $X \in \M_{n\,k} (\F).$ Here $W_{n\,k} \subseteq \M_{n\,k}(\F)$ denotes the space of matrices, all rows of which are equal, being a radical of $\det_{n\,k}$.
\end{theorem}


Our proof of the main theorem relies on the notion of the radical of a function which was introduced and studied by Waterhouse~\cite{Waterhouse1983} (Definition~\ref{def:radical} below). The reduction of the case $(n,k)$ with $k \ge 4, n \ge k + 2$ and $n + k$ odd to the case $(n-1,k)$ with $k \ge 4, n \ge k + 2$ and $(n-1) + k$ even is done by finding the radical of $\det_{n\,k}$ explicitly and using its properties.

Thus, the results of the current paper and its predecessor~\cite{Guterman2025} provide the solution to linear preserver problem for the Cullis' determinant if $k \ge 4, n \ge k + 2$ and the ground field is large enough. The cases $k = 1$ and $k = 2$ are considered in~\cite{Guterman2025} as well. The case $n = k + 1$ is considered in~\cite{Guterman2024}, where the first attempt to find the description of linear maps preserving the Cullis' determinant is made. The remaining case $k = 3$ need a special approach because the proof for the case $k \ge 4$ relies on the following lemma and will be discussed separately. We remark that in this paper we apply the results obtained in~\cite{Guterman2025} and we do not use any results from~\cite{Guterman2024}.

\begin{lemma}[{\cite[Lemma~5.13]{Guterman2025}}]\label{lem:CullisDeg1Rank1}Assume that $|\F| > k \ge 4$, $n \ge k + 2$ and $n + k$ is even. Let  $T\colon \M_{n\, k} (\F) \to \M_{n\, k} (\F)$ be a linear map such that  $\det_{n\, k} (T(X)) = \det_{n\,k}(X)$ for all $X \in \M_{n\, k} (\F).$ Then $\rk (X) = 1$ implies $\rk (T(X)) = 1.$
\end{lemma}

The example below shows that this lemma does not hold for $k = 3$.

\begin{example}[{\cite[Example~5.11]{Guterman2025}}]\label{cex:K3NKEvenInt}Let $B = \begin{psmallmatrix}
                                 1 & 0 & 0 & 0 & -1 & 0 & \cdots & 0\\
                                 0 & 1 & 0 & -1 & 0 & 0 & \cdots & 0\\
                                 0 & 0 & 0 & 0 & 0 & 0 & \cdots & 0
                                        \end{psmallmatrix}^t \in \M_{n\,3}(\F).$ Then $\rk (B) = 2$ and $\deg_{\lambda} (\det_{n\,k}(A + \lambda B)) \le 1$ for all $A \in \M_{n\,3}(\F)$.
\end{example}

Nevertheless, if $k = 3$, then linear maps preserving the Cullis determinant admit the same description as in the case when $k \ge 4$ and will be studied separately.\bigskip

The paper is organized as follows: in Section~\ref{sec:prelim} we provide the basic facts regarding $\det_{n\,k}$; in Section~\ref{sec:kernels} we find the radical of $\det_{n\,k}$ if $n + k$ is odd; in Section~\ref{sec:maintheorem} we reduce the linear preserver problem for $\det_{n\,k}$ and $n + k$ is odd to the linear preserver problem for $\det_{(n-1)\,k}$ and prove the main characterization theorem for $n + k$ odd.

\section{Preliminaries}\label{sec:prelim}

Let us provide the basic notation and definitions used throughout this paper.%\bigskip %% in order to formulate the Main Theorem.

By $\F$ we denote a field without any restrictions on its characteristic. We will use an inequality $|\F| > k$ for $k \in \mathbb N$ which means that either $\F$ is infinite or $\F$ is finite and $|\F| > k$.

We denote by $\M_{n\, k}(\F)$ the set of all $n\times k$ matrices with the entries from a certain field $\F.$ If $X \in \M_{n\,k}(\F)$ then by $x_{i\, j}$ we denote an entry of $X$ lying on the intersection of the $i$-th row and the $j$-th column of $X$.

$O_{n\, k}\in \M_{n\, k}(\F)$ denotes the matrix with all the entries equal to zero. $I_{n\, n} = I_{n} \in \M_{n\, n}(\F)$ denotes an identity matrix. Let us denote by $E_{ij} \in \M_{n\, k}(\mathbb F)$ a matrix, whose entries are all equal to zero besides the entry on the intersection of the $i$-th row and the $j$-th column, which is equal to one. By $x_{i\, j}$ we denote the element of a matrix $X$ lying on the intersection of its $i$-th row and $j$-th column.  For integers $i, j$ we denote by $\delta_{i\,j}$  a \emph{Kronecker delta} of $i$ and $j$, which is equal to $1$ if $i = j$ and equal to $0$ otherwise.

For $A \in \M_{n\,k_1}(\F)$ and $B \in \M_{n\,k_2}(\F)$ by $A|B \in \M_{n\,k_1+k_2}(\F)$ we denote a block matrix defined by $A|B = \begin{pmatrix} A & B\end{pmatrix}$.

For $A \in \M_{n\,k}$ by $A^t \in \M_{k\,n}(\F)$ we denote a transpose of the matrix $A$, i.e. $A^{t}_{i\,j} = A_{j\,i}$ for all $1 \le i \le k$, $1 \le j \le n$.


We use the notation for submatrices following~\cite{Minc1984} and~\cite{Pierce1979}. That is, by $A[J_1|J_2]$ we denote the $|J_1 |\times |J_2|$ submatrix of $A$ lying on the intersection of rows with
the indices from $J_1$ and the columns with the indices from $J_2$. By $A(J_1|J_2)$ we denote a submatrix of $A$ derived from it by striking out the rows with indices belonging to $J_1$ and the columns with the indices belonging to $J_2$. If one of the two index sets is absent, then it means an empty set, i.e. $A(J_1|)$ denotes a matrix derived from $A$ by striking out from it the rows with indices belonging to $J_1$. We may skip curly brackets, i.e. $A[1,2|3,4] = A[\{1,2\}|\{3,4\}]$. The notation with mixed brackets is also used, i.e. $A(|1]$ denotes the first column of the matrix $A$. The above notations are also used for vectors as well. In this case vectors are considered as $n\times 1$ or $1 \times n$ matrices.

To avoid the ambiguity, we would like to clarify that the symbol $\circ$ is used  only in the sense of composition of two maps, that is, if $f \colon A \to B$ and $g \colon B \to C$, then by $g\circ f\colon A \to C$ we denote a composition of $f$ and $g$ defined by $(g\circ f)(x) = g(f(x))$ for all $x \in A$. 

\begin{definition}
By $[n]$ we denote the set $\{1, \ldots, n\}$.
\end{definition}
\begin{definition}
 By $\mathcal C_{X}^{k}$ we denote the set of injections from $[k]$ to $X$.
\end{definition}
\begin{definition}By $\binom{X}{k}$ we denote the set of the images of injections from $[k]$ to $X$, i.e. the set of all subsets of $X$ of cardinality $k.$
\end{definition}


\begin{definition}Suppose that $c \in \binom{[n]}{k}$ equal to $\{i_1,\ldots, i_k\}$, where $i_1 < i_2 < \ldots < i_k$, and $1 \le \alpha \le k$ is a natural number. Then $c(\alpha)$ is defined by
$$c(\alpha) = i_{\alpha}.$$
\end{definition}

\begin{definition}Given a set $c \in \binom{[n]}{k}$ we denote by $\sgn (c) = \sgn_{[n]} (c)$ the number
$$(-1)^{\sum_{\alpha = 1}^{k} (c(\alpha) - \alpha)}.$$
\end{definition}


\begin{note}$\sgn_{[n]}(c)$ depends only on $c$ and does not depend on $n.$
\end{note}


\begin{definition}Given an injection $\sigma \in \mathcal C_{[n]}^{k}$ we denote by $\sgn_{n\,k} (\sigma)$ the product $\sgn(\pi_\sigma) \cdot \sgn_{[n]} (c),$
where $\sgn(\pi)$ is the sign of the permutation
$$
\pi_\sigma =
\begin{pmatrix}
i_1 & \ldots & i_k\\
\sigma(1) & \ldots & \sigma(k)
\end{pmatrix},
$$
where $\{i_1, \ldots, i_k\} = \sigma([k])$ and $i_1<i_2<\ldots <i_k$.
\end{definition}

We omit the subscripts for $\sgn_{[n]}$ and $\sgn_{n\,k}$ if this cannot lead to a misunderstanding.%\bigskip
%\begin{underconst}
%Now we can provide the definition of Cullis determinant following~\cite{NAKAGAMI2007422}.
%\end{underconst}
\begin{definition}[\cite{NAKAGAMI2007422}, Theorem 13]\label{def:CullisDet} Let  $n \ge k$, $X \in \M_{n\,k} (\mathbb F)$. Then Cullis' determinant $\det_{n\, k}(X)$ of $X$ is defined to be the function:
$$
\det_{n\, k} (X) = \sum_{\sigma \in \mathcal C_{k}^{[n]}} \sgn_{n\,k} (\sigma) x_{\sigma(1)\, 1}x_{\sigma(2)\, 2}\ldots x_{\sigma(k)\, k}.
$$
We also denote $\det_{n\, k} (X)$ as follows
\[
\det_{n\,k}(X)=\begin{vmatrix}x_{1\,1} & \cdots & x_{1\,k}\\
\vdots & \cdots & \vdots\\
x_{n\,1} & \cdots & x_{n\,k}
\end{vmatrix}_{n\,k}.
\]
If $n = k$, then we also write $\det_{k}$ or $\det$ instead of $\det_{n\,k}$ because in this case $\det_{n\,k}$ is clearly equal to an ordinary determinant of a square matrix.
\end{definition}



Let us first list the properties of $\det_{n\, k}$ which are similar to corresponding properties of the ordinary determinant (see \cite[\textsection 5, \textsection 27, \textsection 32]{cullis1913} or~\cite{NAKAGAMI2007422} for detailed proofs).


\begin{theorem}[{\cite[Theorem 13, Theorem 16]{NAKAGAMI2007422}}]
\begin{enumerate}
\item[]
\item For $X \in \M_{n}(\mathbb F),$ $\det_{n\,n}(X) = \det (X).$
\item For $X \in \M_{n\,k}(\mathbb F),$ $\det_{n\,k}(X)$ is a linear function of columns of $X$.
\item If a matrix $X \in \M_{n\,k}(\mathbb F)$ has two identical columns or one of its columns is a linear combination of other columns, then $\det_{n\,k}(X)$ is equal to zero.
\item For $X \in \M_{n\,k}(\mathbb F),$ interchanging any two columns of $X$ changes the sign of $\det_{n\,k}(X)$.
\item Adding a linear combination of columns of $X$ to another column of $X$ does not change $\det_{n\,k}(X)$.
\item For $X \in \M_{n\,k}(\mathbb F),$ $\det_{n\,k}(X)$ can be calculated using the Laplace expansion along a column of $X$ (see Lemma~\ref{lem:DetNKLaplaceExp} for precise formulation).
\end{enumerate}
\end{theorem}


\begin{corollary}\label{cor:CullisBinomialExpansion}Let $n \ge k$, $A, B \in \M_{n\,k} (\F).$% Consider $P \in \F[\lambda]$ defined by an equality $P(\lambda) = \det_{n\, k} (A + \lambda B)$.
Then
\begin{multline}\label{eq:CullisBinomialExpansion}
\det_{n\,k} (A + \lambda B)\\
= \sum_{d = 0}^{k}\lambda^d \left( \sum_{1 \le i_1 < \ldots < i_d \le k} \det_{n\,k}\Bigl(A(|1]\Big|\ldots \Big| B(|i_1] \Big| \ldots \Big| B(|i_d] \Big| \ldots \Big| A(|k] \Bigr)\right),
\end{multline}
where both sides of the equality are considered as formal polynomials in $\lambda$, i.e. as elements of $\F[\lambda]$.
\end{corollary}
\begin{proof}This is a direct consequence from the multilinearity of $\det_{n\,k}$ with respect to the columns of a matrix.
\end{proof}
\begin{corollary}\label{cor:DegDetABLEQK}If $A, B \in \M_{n\,k}(\F)$, then $\deg_\lambda\left(\det_{n\,k} (A + \lambda B)\right) \le k$.
\end{corollary}


\begin{lemma}[{\cite[Corollary~2.8]{Guterman2025}}]\label{lem:rightmatrixmult}If $X \in \M_{n\,k}(\F)$ and $Y \in \M_{k\,k}(\F)$, then
\[
\det_{n\,k}(XY) =  \det_{n\,k}(X)\det_k(Y).
\]
\end{lemma}



\begin{lemma}[{Cf.~\cite[Theorem~16]{NAKAGAMI2007422}}]\label{lem:DetNKLaplaceExp}
Let $1 < k \le n$. For any $n \times k$ matrix $X = (x_{i\,j})$ the expansion of $\det_{n\,k}(X)$ along the $j$-th column
is given by
\[
\det_{n\,k}(X) = \sum_{i = 1}^n (-1)^{i+j} x_{i\,j} \det_{(n-1)\,(k-1)}\Bigl(X(i|j)\Bigr).
\]
%where $X_{(i, j)}$ denotes the $(n-1) \times (k-1)$ matrix obtained from $Y$ by deleting the $i$-th row and the $j$-th column.
\end{lemma}





\begin{lemma}[{Invariance of $\det_{n\,k}$ under cyclic shifts, Cf.~\cite[Theorem 3.5]{amiri2010}}]\label{lem:PermKNOddCyclic}If $k \le n$, and $k + n$ is odd, then for all $X = (x_{i\,j}) \in \M_{n\,k}(\F)$ and $i \in \{1, \ldots, n\}$ 
\[
(-1)^{(i+1)k}
\begin{vmatrix}
x_{i\,1} & \ldots & x_{i\,k}\\
\vdots & \vdots & \vdots\\
x_{n\,1} & \ldots & x_{n\,k}\\
x_{1\,1} & \ldots & x_{1\,k}\\
  \vdots & \vdots & \vdots\\
x_{i-1\,1} & \ldots & x_{i-1\,k}\\
\end{vmatrix}_{n\,k} =
\begin{vmatrix}
x_{1\,1} & \ldots & x_{1\,k}\\
  \vdots & \vdots & \vdots\\
x_{n\,1} & \ldots & x_{n\,k}\\
\end{vmatrix}_{n\,k}.
\]
Here the matrix on the left-hand side of the equality is obtained from $X$ by performing the row cyclical shift sending $i$-th row of $X$ to the first row of the result.
\end{lemma}


\begin{lemma}[{Cf.~\cite[Lemma 3.2]{amiri2010}}]\label{lem:AMIRILEMMA32} Let $n \ge k \ge 1$ and $x_{i\,j} \in \F$ for all $1 \le i \le n, 1 \le j \le k$. Then
\[
\begin{vmatrix}
x_{1\,1} & \ldots & x_{1\,k}\\
\vdots & \vdots & \vdots\\
x_{n\,1} & \ldots & x_{n\,k}\\
0 & \ldots & 0
\end{vmatrix}_{n+1\,k} = 
\begin{vmatrix}
x_{1\,1} & \ldots & x_{1\,k}\\
\vdots & \vdots & \vdots\\
x_{n\,1} & \ldots & x_{n\,k}
\end{vmatrix}_{n\,k}.
\]
\end{lemma}


\begin{lemma}[Cf.~{\cite[Lemma~20]{NAKAGAMI2007422}}]~\label{lem:NAKAGAMILEMMA20} Assume that $n > k \ge 1$. Let  $X \in \M_{n\,k}(\F)$ and $Y \in \M_{n\,(k+1)}$ is defined by $Y = X | \begin{psmallmatrix}1 \\ \vdots \\ 1\end{psmallmatrix}$. Then
\[
\det_{n\,(k+1)}(Y) = \begin{cases}
\det_{n\,k}(X), & \mbox{$n + k$ is odd},\\
0, & \mbox{$n + k$ is even}.
\end{cases}
\]
\end{lemma}


\begin{lemma}[{Cf.~\cite[Theorem~3.3]{amiri2010}}]\label{lem:AdditionNPLusKOdd}
{Suppose $1 \le k < n,$ and $k + n$ be an odd integer, $A, X \in \M_{n\, k}(\F)$ and $X = \begin{psmallmatrix}
x_1 & \ldots & x_k\\
\vdots & \ddots & \vdots\\
x_1 & \ldots & x_k
\end{psmallmatrix}$ for some $x_1,\ldots x_k \in \F$. Then}
\[
\det_{n\,k} (A + X) = \det_{n\,k} (A).
\]
\end{lemma}


\begin{corollary}\label{cor:AdditionLinearMapNPLusKOdd}Assume that $n + k$ is odd. Let 
$\phi \colon \M_{n\,k}(\F) \to W_{n\,k}$ and $T \colon \M_{n\,k}(\F) \to \M_{n\,k}(\F)$ be linear maps such that  $\det_{n\, k} (T(X)) = \det_{n\,k}(X)$ for all $X \in \M_{n\,k}(\F)$. Denote by $S$ a linear map $\M_{n}(\F) \to \M_{n\,k}(\F)$ defined by
\begin{equation}\label{eq:AddingKernel}
S(X) = T(X) + \phi (X)\;\;\mbox{for all}\;\;X \in \M_{n\,k}(\F).
\end{equation}
 Then
\[
\det_{n\, k} (S(X)) = \det_{n\,k}(X)\;\;\mbox{for all}\;\;X \in \M_{n\,k} (\F).
\]
\end{corollary}

\begin{lemma}[{\cite[Lemma~3.1]{Guterman2025}}]\label{lem:TwoSidedMulPreservesDet}Let $n \ge k$, $A \in \M_{n\,n}(\F)$, $B \in \M_{k\,k}(\F)$, and let $T\colon \M_{n\,k}(\F) \to \M_{n\,k}(\F)$ be a linear map defined by
\begin{equation}\label{eq:TwoSideMul}
T(X) = AXB
\end{equation}
for all $X \in \M_{n\,k}(\F).$ Then
\[
\det_{n\,k} (T(X)) = \det_{n\,k}(X)
\]
for all $X \in \M_{n\,k} (\F)$ if and only if
\begin{equation}\label{eq:TwoSideMul2}
\det_{n\,k} \Bigl(A(|d]\Bigr)\cdot \det_k \Bigl(B\Bigr) =  \sgn (d)
\end{equation}
for all $d \in \binom{[n]}{k}.$
\end{lemma}

\begin{lemma}[{\cite[Theorem~5.14]{Guterman2025}}]\label{lem:MainTheoremEven}Assume that $|\F| > k \ge 4, n \ge k + 2$ and  $n + k$ is even. Let $T\colon \M_{n\, k} (\F) \to \M_{n\, k} (\F)$ be a linear map. Then $\det_{n\, k} (T(X)) = \det_{n\,k}(X)$ for all $X \in \M_{n\, k} (\F)$ if and only if there exist $A \in \M_{n\, n}(\F)$ and $B \in \M_{k\, k}(\F)$ such that
\begin{equation}\label{thm:MainTheoremEvenKGe4:eq}
\det_{n\, k} \Bigl(A(|i_1,\ldots, i_k]\Bigr) \cdot \det_k \Bigl(B\Bigr) = (-1)^{i_1 + \ldots + i_k - 1 - \ldots - k}
\end{equation}
for all increasing sequences $1 \le i_1 < \ldots < i_k \le n $ and
\[
T(X) = AXB
\]
for all $X \in \M_{n\, k} (\F).$
\end{lemma}


The following general fact regarding polynomials over a field will be useful in the sequel. 
%\begin{lemma}[{Cf.~\cite[IV, \S 1, Corollary 1.8]{Lang2002}}]Let $\F$ be a finite field with $q$ elements. Let $f$ be a
%polynomial in $n$ variables over $\F$ such that the degree of $f$ in each variable
%is less than $q$. If $f$ defines the zero function on $\F^{n}$, then $f = 0$.
%\end{lemma}


\begin{lemma}[{\cite[Corollary~2.13]{Guterman2025}}]\label{DegFGLessEqual}Let $f, g$ be
polynomials in $n$ variables over $\F$ such that either $\F$ is infinite or the degree of $f$ and $g$ in each variable
is less than $|\F|$. If $f$ and $g$ define the same function on $\F^{n}$, then $f = g$.
\end{lemma}

In order to prove that every linear map preserving $\det_{n\,k}$ is invertible we use the notion and the properties of radical of a function on a vector space following~\cite{Waterhouse1983}.

\begin{definition}[Cf.~{\cite[text at the beginning of Section~1]{Waterhouse1983}}]\label{def:radical}Let $\F$ be a field, $V$ be a finite-dimensional vector space over $\F$, and let $f$ be a function from $V$ to $\F$. The \emph{radical} of $f$, denoted by $\rad (f)$ is a subset of $V$ defined by
\[
\rad (f) = \{\mathbf w \mid f(\mathbf v + \lambda \mathbf w) = f(\mathbf v)\;\;\mbox{for all}\;\; \mathbf v \in V,\;\; \lambda \in \F\}.
\]
\end{definition}

\begin{definition}[Cf.~{\cite[text at the beginning of Section~1]{Waterhouse1983}}]\label{def:frad}If $\rad(f)$ is the radical of $f$ and $\pi_{\rad}\colon V \to V / \rad(f)$ is the canonical projection, then $f_{\rad}$ defined as a unique function from $V / \rad(f)$  to $\F$ with trivial radical such that $f = f_{\rad}\circ \pi_{\rad}$.
\end{definition}

\begin{lemma}[Cf.~{\cite[Proposition~1]{Waterhouse1983}}]\label{lem:RadicalCriterion}
$\rad (f)$ is nonzero if and only if there is a noninvertible linear map $T\colon V\to V$ satisfying $f(T\mathbf v) = f(\mathbf v)$ for all $\mathbf v$ in $V$.
\end{lemma}

\begin{lemma}\label{lem:KernelContainedInRad}If $T\colon V \to V$ is a linear map preserving $f$, then $\Ker T \subseteq \rad(f)$.
\end{lemma}
\begin{proof}Indeed, assume that $\mathbf w \in \Ker T.$ Then
\begin{equation*}
f(\mathbf v + \lambda\mathbf w)= f(T(\mathbf v + \lambda\mathbf w))\\
= f(T(\mathbf v) + \lambda T(\mathbf w)) = f(T(\mathbf v) + \mathbf 0) = f(\mathbf v)
\end{equation*}
for all $\mathbf v \in V$ and $\lambda \in \F$. Therefore, $\mathbf w \in \rad(f)$.
\end{proof}

\begin{remark}The lemma above is implicitly proved in~\cite{Waterhouse1983} (see~\cite[Proposition~1]{Waterhouse1983}).
\end{remark}

\begin{lemma}[Cf.~{\cite[Proposition~2]{Waterhouse1983}}]\label{lem:RadicalInvariant}
Let $V$ be a vector space over $\F$, $f\colon V \to \F$ be an arbitrary function. If $f$ has trivial radical, then the set of all linear maps preserving $f$ form a group. If $f$ has nontrivial radical $\rad(f)$, let $\pi_{rad}\colon V \to V/\rad(f)$ be the projection, and write $f = f_{\rad} \circ \pi_{\rad}$. Then the linear map $T \colon V \to V$ preserves $f$ if and only if
\begin{enumerate}[label=(\alph*), ref=(\alph*)]
\item\label{lem:RadicalInvariant:part1} $T(\rad(f)) \subseteq \rad(f)$  and
\item\label{lem:RadicalInvariant:part2} the map $T_{\rad} \colon V / \rad(f) \to V / \rad(f)$ that is given by $T_{\rad}(\pi_{\rad}(\mathbf v)) = \pi_{\rad}(T(\mathbf v))$ (and is well defined because of~\ref{lem:RadicalInvariant:part1}) preserves $f_{\rad}$.
\end{enumerate}
\end{lemma}



\section{The radical of $\det_{n\,k}$ if $n + k$ is odd}\label{sec:kernels}

%In this subsection we assume that $n \ge k$ and $n + k$ is odd. This implies that $n \ge k+1$.

\begin{remark}The approach used in this section is similar to one used in~\cite[Section~4]{Guterman2025} but differs in details.
\end{remark}

\begin{lemma}\label{lem:DetNKXiEqX1}Suppose that $n \ge k \ge 1, n+k$ is odd, $x_1, \ldots, x_n \in \F$ and
\[
X = \begin{pmatrix}
x_1 & 0 & 0 & \cdots & 0\\
x_2 & 0 & 0 & \cdots & 0\\
x_3 & 1 & 0 &\cdots & 1\\
x_4 & 0 & 1 &\cdots & 1\\
\vdots & \vdots & \vdots & \ddots & \vdots\\
x_k & 0 & 0 & \cdots & 1\\
%x_{k+1} & 0 & 0 & \cdots  & 1\\
\vdots & \vdots & \vdots & \ddots & \vdots\\
x_{n} & 0 & 0 & \cdots & 1
\end{pmatrix} \in \M_{n\,k}(\F).
\]
Then
\[
\det_{n\,k}(X) = (-1)^{k-1}(x_1 - x_2).
\]
\end{lemma}
\begin{proof}
Observe that
\begin{equation}\label{lem:DetNKXiEqX1:eq1}
\begin{vmatrix}
x_1 & 0 & 0 & \cdots & 0\\
x_2 & 0 & 0 & \cdots & 0\\
x_3 & 1 & 0 &\cdots & 1\\
x_4 & 0 & 1 &\cdots & 1\\
\vdots & \vdots & \vdots & \ddots & \vdots\\
x_k & 0 & 0 & \cdots & 1\\
%x_{k+1} & 0 & 0 & \cdots  & 1\\
\vdots & \vdots & \vdots & \ddots & \vdots\\
x_{n} & 0 & 0 & \cdots & 1
\end{vmatrix}_{n\,k} =
\begin{vmatrix}
x_1 & 0 & 0 & \cdots & 1\\
x_2 & 0 & 0 & \cdots & 1\\
x_3 & 1 & 0 &\cdots & 1\\
x_4 & 0 & 1 &\cdots & 1\\
\vdots & \vdots & \vdots & \ddots & \vdots\\
x_k & 0 & 0 & \cdots & 1\\
%x_{k+1} & 0 & 0 & \cdots  & 1\\
\vdots & \vdots & \vdots & \ddots & \vdots\\
x_{n} & 0 & 0 & \cdots & 1
\end{vmatrix}_{n\,k}
-
\begin{vmatrix}
x_1 & 0 & 0 & \cdots & 1\\
x_2 & 0 & 0 & \cdots & 1\\
x_3 & 1 & 0 &\cdots & 0\\
x_4 & 0 & 1 &\cdots & 0\\
\vdots & \vdots & \vdots & \ddots & \vdots\\
x_k & 0 & 0 & \cdots & 0\\
%x_{k+1} & 0 & 0 & \cdots  & 0\\
\vdots & \vdots & \vdots & \ddots & \vdots\\
x_{n} & 0 & 0 & \cdots & 0
\end{vmatrix}_{n\,k}.
\end{equation}

Consider the first term in the difference~\eqref{lem:DetNKXiEqX1:eq1}. Lemma~\ref{lem:NAKAGAMILEMMA20} implies that it is zero since $n - k - 1$ is even. 

Now consider the second term in the difference~\eqref{lem:DetNKXiEqX1:eq1}. Using the Laplace expansion along the second column for $k-2$ times we obtain that
\begin{equation}\label{lem:DetNKXiEqX1:eq3}
\begin{vmatrix}
x_1 & 0 & 0 & \cdots & 1\\
x_2 & 0 & 0 & \cdots & 1\\
x_3 & 1 & 0 &\cdots & 0\\
x_4 & 0 & 1 &\cdots & 0\\
\vdots & \vdots & \vdots & \ddots & \vdots\\
x_k & 0 & 0 & \cdots & 0\\
%x_{k+1} & 0 & 0 & \cdots  & 0\\
\vdots & \vdots & \vdots & \ddots & \vdots\\
x_{n} & 0 & 0 & \cdots & 0
\end{vmatrix}_{n\,k} = (-1)^{s} \begin{vmatrix}
x_1 & 1\\
x_2 & 1\\
x_{k+1} & 0\\
%x_{k+1}\\
\vdots & \vdots \\
x_{n} & 0
\end{vmatrix}_{n-k+2\,2},
\end{equation}
where 
\[
s = {(2+3) + \ldots + (2+3)} = (2+3)(k-2).
\]
Hence, \begin{equation}\label{lem:DetNKXiEqX1:eq4}
(-1)^{s} = (-1)^{k-2}.
\end{equation}

The Laplace expansion along the second column yields that
\begin{multline}\label{lem:DetNKXiEqX1:eq2}
\begin{vmatrix}
x_1 & 1\\
x_2 & 1\\
x_{k+1} & 0\\
%x_{k+1}\\
\vdots & \vdots\\
x_{n} & 0
\end{vmatrix}_{n-k+1\,2} = (-1)^{2 + 1}\begin{vmatrix}
x_2 \\
x_{k+1}\\
%x_{k+1}\\
\vdots \\
x_{n} 
\end{vmatrix}_{n-k+1\,1}
+ (-1)^{2 + 2}
\begin{vmatrix}
x_1 \\
x_{k+1}\\
%x_{k+1}\\
\vdots \\
x_{n} 
\end{vmatrix}_{n-k+1\,1}\\ = -\left(x_2 + \sum_{i=1}^{n-k} (-1)^{i} x_{k+i}\right) + \left(x_1 + \sum_{i=1}^{n-k} (-1)^{i} x_{k+i}\right) = x_1 - x_2.
\end{multline}

Thus, after making the substitution of~\eqref{lem:DetNKXiEqX1:eq2}, \eqref{lem:DetNKXiEqX1:eq3} and~\eqref{lem:DetNKXiEqX1:eq4} into~\eqref{lem:DetNKXiEqX1:eq1}, we  obtain that
\[
\det_{n\,k}(X) = 0 - (-1)^{k-2}(x_1 - x_2) = (-1)^{k-1}(x_1 - x_2).
\]
\end{proof}

\begin{lemma}\label{lem:NKEvenAdditionPreserverHasZero11Odd}Let $n \ge k, n + k$ be even, $|\F| > k$ and $Y \in \M_{n\, k} (\F)$ be such that
\begin{equation}\label{lem:NKEvenAdditionPreserverHasZero11Odd:cond}
\det_{n\, k} (A + \lambda Y) = \det_{n\, k} (A)
\end{equation}
for all $A \in \M_{n\, k}(\F)$ and $\lambda \in \F$. Then $y_{1\,1} - y_{2\,1} = 0$.
\end{lemma}
\begin{proof}Let $A \in \M_{n\, k}(\F)$ be defined by
\[
A = E_{3\,2} + E_{4\,3} + \ldots + E_{k\, k-1} + E_{3\, k} + \ldots + E_{n\, k}  = \begin{pmatrix}
0 & 0 & 0 & \cdots & 0\\
0 & 0 & 0 & \cdots & 0\\
0 & 1 & 0 &\cdots & 1\\
0 & 0 & 1 &\cdots & 1\\
\vdots & \vdots & \vdots & \ddots & \vdots\\
0 & 0 & 0 & \cdots & 1\\
%0 & 0 & 0 & \cdots  & 1\\
\vdots & \vdots & \vdots & \ddots & \vdots\\
0 & 0 & 0 & \cdots & 1
\end{pmatrix}.
\]

It follows from Corollary~\ref{cor:CullisBinomialExpansion} that there is a polynomial $P \in \F[\lambda]$ with $\deg_{\lambda}(P) \le k$ such that $P = \det_{n\, k} (A + \lambda Y)$. Let $a_0, \ldots, a_k \in \F$ be the coefficients of $P$. That is,
 \[P =   \det_{n\, k} (A + \lambda Y) = a_0 + a_1\lambda + \ldots + a_k \lambda^k.
  \]
On the one hand, the condition~\eqref{lem:NKEvenAdditionPreserverHasZero11Odd:cond} implies that
\[
a_0 + a_1\lambda + \ldots + a_k \lambda^k = \det_{n\,k} (A)
\]
for all $\lambda \in \F$. Therefore, 
\begin{equation}\label{lem:NKEvenAdditionPreserverHasZero11Odd:eq3}
a_1 = \ldots = a_k = 0
\end{equation}
 by Lemma~\ref{DegFGLessEqual} because  $|\F| > k$.

On the other hand, the definition of $a_1$ implies that
\[
a_1 = \sum_{1 \le i_1 \le k} \det_{n\,k}\Bigl(A(|1]\Big|\ldots \Big| Y(|i_1] \Big| \ldots \Big| A(|k] \Bigr)
\]
by~\eqref{eq:CullisBinomialExpansion}. Since $A(|1]$ is a zero column, this implies that 
\begin{multline*}
a_1 = \sum_{1 \le i_1 \le k} \det_{n\,k}\Bigl(A(|1]\Big|\ldots \Big| Y(|i_1] \Big| \ldots \Big| A(|k] \Bigr)\\
 = \det_{n\,k} \Bigl(Y(|1] \Big| A(|2] \Big| \ldots \Big| A(|k]\Bigr)  + \sum_{2 \le i_1 \le k} \det_{n\,k}\Bigl(A(|1]\Big|\ldots \Big| Y(|i_1] \Big| \ldots \Big| A(|k] \Bigr)\\
 = \det_{n\,k} \Bigl(Y(|1] \Big| A(|2] \Big| \ldots \Big| A(|k]\Bigr)  + \sum_{2 \le i_1 \le k} \det_{n\,k}\Bigl(0\Big|\ldots \Big| Y(|i_1] \Big| \ldots \Big| A(|k] \Bigr)\\
 =  \det_{n\,k} \Bigl(Y(|1] \Big| A(|2] \Big| \ldots \Big| A(|k]\Bigr) + \sum_{2 \le i_1 \le k} 0\\
 = \det_{n\,k} \Bigl(Y(|1] \Big| A(|2] \Big| \ldots \Big| A(|k]\Bigr).
\end{multline*}

By substituting of the right-hand side of the above equality in~\eqref{lem:NKEvenAdditionPreserverHasZero11Odd:eq3} we obtain that
\begin{equation}\label{lem:NKEvenAdditionPreserverHasZero11Odd:eq1}
\det_{n\,k} \Bigl(Y(|1] \Big| A(|2] \Big| \ldots \Big| A(|k]\Bigr) = 0.
\end{equation}

Now note that $\Bigl(Y(|1] \Big| A(|2] \Big| \ldots \Big| A(|k]\Bigr)$ has the form described in the statement of Lemma~\ref{lem:DetNKXiEqX1} for $x_1 = y_{1\, 1}, \ldots, x_n = y_{n\,1}$. This implies that\begin{equation}\label{lem:NKEvenAdditionPreserverHasZero11Odd:eq2}
\det_{n\, k} \Bigl(Y(|1] \Big| A(|2] \Big| \ldots \Big| A(|k]\Bigr) = (-1)^{k-1}(y_{1\,1} - y_{2\,1}).
\end{equation}

The equalities~\eqref{lem:NKEvenAdditionPreserverHasZero11Odd:eq1} and~\eqref{lem:NKEvenAdditionPreserverHasZero11Odd:eq2} allow us to conclude that
\[
(-1)^{k-1}(y_{1\,1} - y_{2\,1}) = \det_{n\, k} \Bigl(Y(|1] \Big| A(|2] \Big| \ldots \Big| A(|k]\Bigr) = 0
\]
and consequently
\[
y_{1\,1} - y_{2\,1} = 0.
\]
\end{proof}

\begin{definition}\label{def:SCSDef}Let $n \ge k$ and $1 \le i \le n$, $1 \le j \le k$. By $\SCS_{i\,j}$ we denote a linear map on $\M_{n\,k}(\F)$ defined by 
\begin{multline*}
\SCS_{i\,j}\begin{pmatrix}
x_{1\,1} & \cdots  & x_{1\,k}\\
  \vdots & \ddots & \vdots\\
x_{n\,1} & \cdots & x_{n\,k}\\
\end{pmatrix}\\ = (-1)^{i+1}\begin{pmatrix}
(-1)^{1 - \delta_{1\,j}}x_{i\, j} & \ldots & x_{i\,1} & \ldots & x_{i\, k}\\
\vdots & \ddots & \vdots & \ddots  & \vdots\\
(-1)^{1 - \delta_{1\,j}}x_{n\, j} & \ldots & x_{n\,1} & \ldots & x_{n\, k}\\
(-1)^{1 - \delta_{1\,j}}x_{1\, j} & \ldots &  x_{1\,1} & \ldots & x_{1\, k}\\
\vdots & \ddots & \vdots & \ddots & \vdots\\
(-1)^{1 - \delta_{1\,j}}x_{i-1\, j} & \ldots & x_{i-1\,1} & \ldots &  x_{i-1\, k}
\end{pmatrix}.
\end{multline*}
That is, $\SCS_{i\,j}(X)$ is obtained from $X$ by performing the following sequence of operations:
\begin{enumerate}
\item\label{lst:SIJSeq:it1}the row cyclical shift sending $i$-th row of $X$ to the first row of the result;
\item exchanging the first and the $j$-th column;
\item\label{lst:SIJSeq:it3} multiplying the first column  by $(-1)^{1 - \delta_{1\,j}}$;
\item multiplying all the entries by $(-1)^{i+1}$.
\end{enumerate}
\end{definition}

\begin{lemma}\label{lem:ReduceTo11ByShiftOdd}Assume that $n \ge k$ and $n + k$ is odd. Then $\SCS_{i\,j}$ is an invertible linear map preserving $\det_{n\,k}$ for all $1 \le i \le n$, $1 \le j \le k$.
\end{lemma}
\begin{proof}Invertibility of $\SCS_{i\,j}$ follows directly from the definition. Assume that $X = (x_{i\,j}) \in \M_{n\,k}(\F)$. Let $X' \in \M_{n\,k}(\F)$ be defined by 
\[
X' = 
\begin{pmatrix}
-x_{i\, j} & \ldots & x_{i\,1} & \ldots & x_{i\, k}\\
\vdots & \ddots & \vdots & \ddots  & \vdots\\
-x_{n\, j} & \ldots & x_{n\,1} & \ldots & x_{n\, k}\\
-x_{1\, j} & \ldots &  x_{1\,1} & \ldots & x_{1\, k}\\
\vdots & \ddots & \vdots & \ddots & \vdots\\
-x_{i-1\, j} & \ldots & x_{i-1\,1} & \ldots &  x_{i-1\, k}
\end{pmatrix},
\]
That is, $X'$ is obtained from $X$ by performing operations \ref{lst:SIJSeq:it1}--\ref{lst:SIJSeq:it3} from the sequence in the definition of $\SCS_{i\,j}$ (Definition~\ref{def:SCSDef}).

Let $C_{S} \in \M_{k\,k}(\F)$ be a matrix such that a linear map $Y \mapsto YC_S$ on $\M_{n\,k}(\F)$ exchanges the first and the $j$-th row of the matrix and multiplies its first column by $(-1)^{1 - \delta_{1\,j}}$. That is, $C_S$ is the permutation matrix corresponding to a transposition/trivial permutation $(1j)$ which first row is multiplied by $(-1)^{1 - \delta_{1\,j}}$. 

Then $X' = (-1)^{i+1}\SCS_{i\,j}(X)$ by the definition of $\SCS_{i\,j}$. In addition, $\det_k(C_S) = 1$ by the definition of $C_S$. 

On the one hand, we have the following sequence of equalities
\begin{equation}\label{lem:ReduceTo11ByShiftOdd:eq1}
\det_{n\,k}(\SCS_{i\,j}(X)) = \det_{n\,k}((-1)^{i+1}X') = (-1)^{(i+1)k}\det_{n\,k}(X')
\end{equation}
where the last equality follows from the multilinearity of $\det_{n\,k}$ with respect to columns of $X'$.

On the other hand,  
\[
XC_S =  \begin{pmatrix}
(-1)^{1 - \delta_{1\,j}}x_{1\, j} & \ldots & x_{1\,1} & \ldots & x_{1\, k}\\
\vdots & \ddots & \vdots & \ddots  & \vdots\\
(-1)^{1 - \delta_{1\,j}}x_{(n-i+1)\, j} & \ldots & x_{(n-i+1)\,1} & \ldots & x_{(n-i+1)\, k}\\
(-1)^{1 - \delta_{1\,j}}x_{(n-i+2)\, j} & \ldots &  x_{1\,1} & \ldots & x_{1\, k}\\
\vdots & \ddots & \vdots & \ddots & \vdots\\
(-1)^{1 - \delta_{1\,j}}x_{n\, j} & \ldots & x_{i-1\,1} & \ldots &  x_{i-1\, k}
\end{pmatrix}.
\]
by the definition of $C_S$. Hence, \begin{multline}\label{lem:ReduceTo11ByShiftOdd:eq2}
(-1)^{(i+1)k}\det_{n\,k}(X') = (-1)^{(i+1)k}\begin{vmatrix}
-x_{i\, j} & \ldots & x_{i\,1} & \ldots & x_{i\, k}\\
\vdots & \ddots & \vdots & \ddots  & \vdots\\
-x_{n\, j} & \ldots & x_{n\,1} & \ldots & x_{n\, k}\\
-x_{1\, j} & \ldots &  x_{1\,1} & \ldots & x_{1\, k}\\
\vdots & \ddots & \vdots & \ddots & \vdots\\
-x_{i-1\, j} & \ldots & x_{i-1\,1} & \ldots &  x_{i-1\, k}
\end{vmatrix}\\
 = \det_{n\,k}(XC_S),
\end{multline}
where the last equality follows from Lemma~\ref{lem:PermKNOddCyclic}. In addition, since $\det_k(C_S) = 1$, then
\begin{equation}\label{lem:ReduceTo11ByShiftOdd:eq3}
\det_{n\,k}(XC_S) = \det_{n\,k}(X)
\end{equation}
by Lemma~\ref{lem:rightmatrixmult}.

Therefore, \[
\det_{n\,k}(\SCS_{i\,j}(X)) = (-1)^{(i+1)k}\det_{n\,k}(X') = \det_{n\,k}(XC_S) = \det_{n\,k}(X)
\]
by \eqref{lem:ReduceTo11ByShiftOdd:eq1} and \eqref{lem:ReduceTo11ByShiftOdd:eq2}, \eqref{lem:ReduceTo11ByShiftOdd:eq3}. From this we conclude that $\SCS_{i\,j}$ preserves $\det_{n\,k}$.
\end{proof}

\begin{definition}\label{def:WNK}By $W_{n\,k}\subseteq \M_{n\,k}(\F)$ we denote $k$-dimensional vector space consisting of matrices, all rows of which are equal. That is,
\[
W_{n\,k} = \{\begin{pmatrix}y_1 & \cdots & y_k\\ \vdots & \ddots & \vdots\\ y_1 & \cdots & y_k\end{pmatrix} \mid y_1, \ldots, y_k \in \F\}.
 \]
\end{definition}

\begin{lemma}\label{lem:NKOddRadDetNK}Assume that $|\F| > k$, $n \ge k$ and $n + k$ is odd. Then 
\begin{equation}\label{lem:NKOddRadDetNK:eqq}
\rad(\det_{n\,k}) =  W_{n\,k}.
\end{equation}
\end{lemma}
\begin{proof}
%Let $W_{n\,k} \subseteq \M_{n\,k}(\F)$ be a vector space defined by
%\[
%W_{n\,k} = \{\begin{pmatrix}y_1 & \cdots & y_k\\ \vdots & \ddots & \vdots\\ y_1 & \cdots & y_k\end{pmatrix} \mid y_1, \ldots, y_k \in \F\}.
%\]
The inclusion $W_{n\,k} \subseteq \rad(\det_{n\,k})$  follows from Lemma~\ref{lem:AdditionNPLusKOdd}. Let us prove the inclusion 
\begin{equation}\label{lem:NKOddRadDetNK:eq1}
\rad(\det_{n\,k}) \subseteq  W_{n\,k}.
\end{equation}
For this, assume that $Y \in \rad(\det_{n\,k})$ and show that $y_{i\,j} - y_{i+1\,j}= 0$ for all $1 \le i < n, 1 \le j \le k.$ This clearly implies~\eqref{lem:NKOddRadDetNK:eq1}.

Let $1 \le i \le n$ and $1 \le j \le k$. Consider a matrix $\SCS_{i\,j}(Y)$, where $\SCS_{i\,j}$ is a linear map defined in Definition~\ref{def:SCSDef}. This definition implies that

\begin{equation}\label{lem:NKEvenAdditionPreserverIsZeroOdd:eq1}
\left(\SCS_{i\,j}(Y)\right)_{1\,1} = (-1)^{i - \delta_{1\,j}}y_{i\,j}\quad\mbox{and}\quad \left(\SCS_{i\,j}(Y)\right)_{2\,1} = (-1)^{i - \delta_{1\,j}}y_{(i+1)\,j}.
\end{equation}
Also
\begin{multline*}
\det_{n\,k} (A + \lambda \SCS_{i\,j}(Y)) = \det_{n\,k} (\SCS_{i\,j}(\SCS_{i\,j}^{-1}(A)) + \lambda \SCS_{i\,j}(Y))\\
= \det_{n\,k} (\SCS_{i\,j}(\SCS_{i\,j}^{-1}(A) + \lambda Y)) = \det_{n\,k} (\SCS_{i\,j}^{-1}(A) + \lambda Y)\\
= \det_{n\,k} (\SCS_{i\,j}^{-1}(A)) = \det_{n\,k}(A)
\end{multline*}
for all $A \in \M_{n\,k}(\F)$ and $\lambda \in \F$.

Therefore, $\SCS_{i\,j}(Y)$ satisfies the conditions of Lemma~\ref{lem:NKEvenAdditionPreserverHasZero11Odd}. It implies that
\begin{equation}\label{lem:NKEvenAdditionPreserverIsZeroOdd:eq2}
\left(\SCS_{i\,j}(Y)\right)_{1\,1} - \left(\SCS_{i\,j}(Y)\right)_{2\,1}  = 0.
\end{equation}
By substituting~\eqref{lem:NKEvenAdditionPreserverIsZeroOdd:eq1} to \eqref{lem:NKEvenAdditionPreserverIsZeroOdd:eq2} we obtain that
\[
(-1)^{i - \delta_{1\,j}}\left(y_{i\,j} - y_{(i+1)\,j}\right) = \left(\SCS_{i\,j}(Y)\right)_{1\,1} - \left(\SCS_{i\,j}(Y)\right)_{2\,1}= 0
\]
and consequently
\[
y_{i\,j} - y_{(i+1)\,j} = 0.
\]
This equality holds for all  $1 \le i \le n$ and $1 \le j \le k$. Therefore, $Y \in W_{n\,k}$, and consequently the inclusion~\eqref{lem:NKOddRadDetNK:eq1} holds.
\end{proof}
%
%\begin{lemma}\label{lem:NKOddKernel}Assume that $|\F| > k$, $n \ge k$ and $n + k$ is odd. If $T \colon \M_{n\,k}(\F)$ is a linear map such that $\det_{n\,k}(T(X)) = \det_{n\,k}(X)$, then $\ker (T) \subseteq W_{n\,k}$.
%\end{lemma}
%\begin{proof}Indeed, since $T$ preserves $\det_{n\,k}$,
%\begin{equation}\label{lem:NKOddKernel:eq1}
%\ker(T) \subseteq \rad \left(\det_{n\,k}\right)
%\end{equation}
%by Lemma~\ref{lem:KernelContainedInRad}. Therefore,
%\[
%\ker(T) \overset{\eqref{lem:NKOddKernel:eq1}}{\subseteq} \rad \left(\det_{n\,k}\right) \overset{\eqref{lem:NKOddRadDetNK:eqq}}{=\joinrel=\joinrel=} W_{n\,k}.
%\]
%\end{proof}

\section{The proof of the main theorem}\label{sec:maintheorem}
In this section we describe a relationship between linear $\det_{n\,k}$-preservers and $\det_{(n-1)\,k}$-preservers. This relationship allows us to reduce the linear preserver problem for $\det_{n\,k}$ to the linear preserver problem for $\det_{(n-1)\,k}$ as it is done in Lemma~\ref{lem:FromNKOddToNKEvenIfKerW}. In conclusion we obtain a proof of the main theorem of this paper.

The relationship is established through the space $\M^0_{n\,k}(\F)$ defined below which is possible to identify with $\M_{(n-1)\,k}(\F)$ by the linear maps converting $\det_{n\,k}$ to $\det_{n-1\,k}$ and vice versa. These linear maps are defined in Definition~\ref{def:Lminus} and Definition~\ref{def:Lplus}, respectively.

\begin{definition}
  By $\M^{0}_{n\,k}(\F) \subseteq \M_{n\,k}(\F)$ we denote a linear space of all matrices such that their last row equal to zero, i.e.
  \[
    \M^{0}_{n\,k}(\F) = \{X \in \M_{n\,k}(\F) \mid x_{n\,j} = 0, 1 \le j \le k\}.
  \]
\end{definition}

\begin{definition}\label{def:Lminus}
For every matrix $X \in \M_{n\,k}(\mathbb F)$ by $L^{-}(X) \in \M_{n-1\,k} (\mathbb F)$ we denote the matrix defined by
\[
L^{-}(X) = \begin{pmatrix}X[1|) - X[n|)\\ \vdots \\X[n-1|) - X[n|)\end{pmatrix}.
\]

\end{definition}

\begin{lemma}\label{claim:LminusMatrixForm}$L^{-}(X) = M^{-}X$ for all $X \in \M_{n\,k}(\F)$, where
\[
M^{-} =
\begin{pmatrix}
1 & 0 & \ldots & 0 & -1\\
0 & 1 & \ldots & 0 & -1\\
\vdots & \vdots & \ddots & \vdots & \vdots\\
0 & 0 & \ldots & 1 & -1
\end{pmatrix}
\in \M_{n-1\,n}(\F).
\]
  \end{lemma}
\begin{proof}It could be shown by direct computation. Indeed, if $X = (x_{i\,j}) \in \M_{n\,k}(\F),$ then
\begin{multline*}
M^{-}X = \begin{pmatrix}
1 & 0 & \ldots & 0 & -1\\
0 & 1 & \ldots & 0 & -1\\
\vdots & \vdots & \ddots & \vdots & \vdots\\
0 & 0 & \ldots & 1 & -1
\end{pmatrix}
\begin{pmatrix}
x_{1\,1} & \cdots & x_{1\,k}\\
\vdots & \ddots & \vdots\\
x_{(n-1)\,1} & \cdots & x_{(n-1)\,k}\\
x_{n\,1} & \cdots & x_{n\,k}\\
\end{pmatrix}\\ = \begin{pmatrix}
x_{1\,1} - x_{n\,1} & \cdots & x_{1\,k} - x_{n\,k}\\
\vdots & \ddots & \vdots\\
x_{(n-1)\,1} - x_{n\,1} & \cdots & x_{(n-1)\,k} - x_{n\,k}\\
\end{pmatrix} = \begin{pmatrix}X[1|) - X[n|)\\ \vdots \\X[n-1|) - X[n|)\end{pmatrix}\\
 = L^{-}(X).
\end{multline*}
\end{proof}


\begin{lemma}\label{lem:DetNKLMinus}Assume that $n \ge k \ge 1$ and $n + k$ is odd. Then $\det_{n\,k}(X) = \det_{(n-1)\,k}(L^{-}(X))$ for all $X \in \M_{n\,k}(\F)$.
  \end{lemma}
  \begin{proof}Lemma~\ref{lem:AdditionNPLusKOdd} implies that $\det_{n\,k} (X) = \det_{n\,k}(X')$, where
  \[
  X' = \begin{pmatrix}X[1|) - X[n|)\\ \vdots \\X[n-1|) - X[n|)\\ 0\end{pmatrix}.
  \]
  Then $\det_{n\,k} (X') = \det_{n-1\,k} (L^{-}(X))$ by Lemma~\ref{lem:AMIRILEMMA32} because the last row of $X'$ is zero and $X'[1,\ldots, n-1|) = L^{-}(X)$.
\end{proof}

\begin{definition}\label{def:Lplus}
For every matrix $X \in \M_{n-1\,k}(\mathbb F)$ by $L^{+}(X) \in \M^0_{n\,k} (\mathbb F)$ we denote the matrix obtained from $X$ by adjoining to it the zero row as the last row; that is
\[
L^{+}(X) = \begin{pmatrix}X\\ 0
\end{pmatrix}.
\]
\end{definition}

\begin{lemma}\label{claim:LplusMatrixForm}$L^{+}(X) = M^{+}X$ for all $X \in \M_{n-1\,k}(\F)$, where
\[
M^{+} =
\begin{pmatrix}
1 & 0 & \ldots & 0\\
0 & 1 & \ldots & 0\\
\vdots & \vdots & \ddots & \vdots\\
0 & 0 & \ldots & 1\\
0 & 0 & \ldots & 0
\end{pmatrix}
\in \M_{n\,n-1}(\F).
\]
  \end{lemma}
  \begin{proof}The proof could be done by the direct computation and similar to the proof of Lemma~\ref{claim:LminusMatrixForm}.
  \end{proof}

\begin{lemma}\label{lem:DetNKLPlus}$\det_{(n-1)\,k}(Y) = \det_{n\,k}(L^{+}(Y))$ for all $Y \in \M_{n-1\,k}(\F)$.
  \end{lemma}
\begin{proof}
Let us note first that on the one hand, $L^{-}(L^{+}(Y)) = Y$ because the last row of $L^{+}(Y)$ is zero. On the other hand,
\[
\det_{n\,k}(L^{+}(Y)) = \det_{n-1\,k}(L^{-}(L^{+}(Y)))
\]
by Lemma~\ref{lem:DetNKLMinus}. Therefore, \[
\det_{(n-1)\,k}(Y) = \det_{(n-1)\,k}(L^{-}(L^{+}(Y))) = \det_{n\,k}(L^{+}(Y)).
\]
\end{proof}

Maps $L^{+}$ and $L^{-}$  provide a correspondence between linear maps preserving $\det_{n\,k}$ and linear map preserving $\det_{(n-1)\,k}$ as it is stated in the next two lemmas.


\begin{lemma}\label{lem:LiftDownDetPres}Assume that  $|\F| > k$, $n > k$ and $n + k$ is odd. Let $T$ be a linear map on $\M_{n\, k} (\F)$ such that $\det_{n\, k} (T(X)) = \det_{n\,k}(X)$ for all $X \in \M_{n\,k}(\F)$. If a linear map $S$ on $\M_{(n-1)\,k}(\F)$ is defined by 
$$S = L^{-} \circ T \circ L^{+},$$
then $\det_{(n-1)\, k} (S(Y)) = \det_{(n-1)\,k}(Y)$ for all $Y \in \M_{(n-1)\,k}(\F)$.
\end{lemma}
\begin{proof}Let $Y \in \M_{(n-1)\,k}(\F)$. Then
\begin{equation}\label{lem:LiftDownDetPres:eq1}
\det_{n\,k}(L^{+}(Y)) = \det_{(n-1)\,k}(Y)
\end{equation}
 by Lemma~\ref{lem:DetNKLPlus}. Also %because $L^{-}(L^{+}(Y)) = Y$. Next,
\begin{equation}\label{lem:LiftDownDetPres:eq2}
\det_{n\,k}(T(L^{+}(Y))) = \det_{n\,k}(L^{+}(Y))
\end{equation}
because $T$ preserves $\det_{n\,k}$.

Lemma~\ref{lem:DetNKLMinus} implies that
\begin{equation}\label{lem:LiftDownDetPres:eq3}
\det_{n\,k}(T(L^{+}(Y))) = \det_{(n-1)\,k}(L^{-}(T(L^{+}(Y))).
\end{equation} The equalities~\eqref{lem:LiftDownDetPres:eq1}, \eqref{lem:LiftDownDetPres:eq2} and~\eqref{lem:LiftDownDetPres:eq3} aligned together imply that
\[
\det_{n\,k}(Y) = \det_{n\,k}(L^{-}(T(L^{+}(Y)))\;\;\mbox{for all}\;\; Y \in \M_{(n-1)\,k}(\F).
\]
Since $ L^{-} \circ T \circ L^{+} = S$ by the definition of $S$, we conclude that $S$ preserves $\det_{(n-1)\,k}$.
\end{proof}

\begin{lemma}\label{lem:LiftUpDetPres}Assume that  $|\F| > k$, $n > k$ and $n + k$ is odd. Let $S$ be a linear map on $\M_{(n-1)\, k} (\F)$ such that $\det_{(n-1)\, k} (S(Y)) = \det_{(n-1)\,k}(X)$ for all $Y \in \M_{(n-1)\,k}(\F)$. If a linear map $T$ on $\M_{n\,k}(\F)$ is defined by $$T = L^{+} \circ S \circ L^{+},$$ then $\det_{n\, k} (T(X)) = \det_{n\,k}(X)$ for all $X \in \M_{n\,k}(\F)$
\end{lemma}
\begin{proof}The proof is similar to the proof of the previous lemma. Let $X \in \M_{n\,k}(\F)$. Then
\begin{equation}\label{lem:LiftUpDetPres:eq1}
\det_{(n-1)\,k}(L^{-}(X)) = \det_{n\,k}(X)
\end{equation}
 by Lemma~\ref{lem:DetNKLMinus}. Also %because $L^{-}(L^{+}(Y)) = Y$. Next,
\begin{equation}\label{lem:LiftUpDetPres:eq2}
\det_{(n-1)\,k}(S(L^{-}(X))) = \det_{(n-1)\,k}(L^{-}(X))
\end{equation}
because $S$ preserves $\det_{(n-1)\,k}$.

Lemma~\ref{lem:DetNKLPlus} implies that
\begin{equation}\label{lem:LiftUpDetPres:eq3}
\det_{(n-1)\,k}(S(L^{-}(X))) = \det_{n\,k}(L^{+}(S(L^{-}(X))).
\end{equation} The equalities~\eqref{lem:LiftUpDetPres:eq1}, \eqref{lem:LiftUpDetPres:eq2} and~\eqref{lem:LiftUpDetPres:eq3} aligned together imply that
\[
\det_{n\,k}(X) = \det_{n\,k}(L^{+}(S(L^{-}(X)))\;\;\mbox{for all}\;\; X \in \M_{n\,k}(\F)
\]
Since $ L^{+} \circ S \circ L^{+} = T$ by the definition of $T$, we conclude that $T$ preserves $\det_{n\,k}$.
\end{proof}



\begin{lemma}\label{lem:FromNKOddToNKEvenIfKerW}Assume that  $|\F| > k$, $n > k$ and $n + k$ is odd. Let $T\colon \M_{n\, k} (\F) \to \M_{n\, k} (\F)$ be a linear map such that $\det_{n\, k} (T(X)) = \det_{n\,k}(X)$ and $\Img(T) \subseteq \M^0_{n\,k}(\F)$. Then the following statements hold.
\begin{enumerate}[label=(\alph*), ref=(\alph*)]
\item\label{lem:FromNKOddToNKEvenIfKerW:part4} $\Ker(T) = \rad(\det_{n\,k})$.
\item\label{lem:FromNKOddToNKEvenIfKerW:part5} $L^{+}\circ L^{-} \circ T = T$.
\item\label{lem:FromNKOddToNKEvenIfKerW:part6} $T \circ L^{+} \circ L^{-} = T$.
\item\label{lem:FromNKOddToNKEvenIfKerW:part2} If a linear map $S$ is defined by 
\begin{equation}\label{lem:FromNKOddToNKEvenIfKerW:eq}
S = L^{-} \circ T \circ L^{+},
\end{equation}
then $T = L^{+} \circ S \circ L^{-}$.
\end{enumerate}
\end{lemma}
\begin{proof}~\paragraph{\ref{lem:FromNKOddToNKEvenIfKerW:part4}} The obvious equality $\M^{0}_{n\,k}(\F) \cap \rad(\det_{n\,k}) = \{0\}$ implies that\linebreak $T(\rad(\det_{n\,k})) = \{0\}$ because 
$$T(\rad(\det_{n\,k})) \subseteq \Img(T) \subseteq \M^{0}_{n\,k}(\F)$$
 and $T(\rad(\det_{n\,k})) \subseteq \rad(\det_{n\,k})$ by Lemma~\ref{lem:RadicalInvariant}\ref{lem:RadicalInvariant:part1}.

Therefore, $\Ker(T) \supseteq \rad(\det_{n\,k})$. Since $\Ker(T) \subseteq \rad(\det_{n\,k})$ by Lemma~\ref{lem:KernelContainedInRad}, then we conclude that $\Ker(T) = \rad(\det_{n\,k})$.

\paragraph{\ref{lem:FromNKOddToNKEvenIfKerW:part5}} % follows from the fact that $\Img(T) = \M^0_{n\,k}$. 
Indeed, if $X \in \M_{n\,k}(\F)$, then
\[
T(X) = \begin{pmatrix}(T(X))[1|)\\ \vdots \\(T(X))[n-1|)\\ 0\end{pmatrix}
\]
because $\Img(T) = \M^0_{n\,k}(\F)$. Hence, \begin{multline*}
L^{+}\circ L^{-} \circ T (X) = L^{+}(L^{-}(T(X)))
 = L^{+}(L^{-}(\begin{pmatrix}(T(X))[1|)\\ \vdots \\(T(X))[n-1|)\\ 0\end{pmatrix}))\\
 = L^{+}(\begin{pmatrix}(T(X))[1|)\\ \vdots \\(T(X))[n-1|)\end{pmatrix})
  = \begin{pmatrix}(T(X))[1|)\\ \vdots \\(T(X))[n-1|)\\ 0\end{pmatrix} = T(X).
\end{multline*}
for all $X \in \M_{n\,k}(\F)$.

\paragraph{\ref{lem:FromNKOddToNKEvenIfKerW:part6}} On the one hand,
\begin{multline*}
T \circ L^{+} \circ L^{-} (X) = T(L^{+}(L^{-}(X)))\\
 = T(L^{+}(\begin{pmatrix}X[1|) - X[n|)\\ \vdots \\X[n-1|) - X[n|)\end{pmatrix})) = T(\begin{pmatrix}X[1|) - X[n|)\\ \vdots \\X[n-1|) - X[n|)\\ 0\end{pmatrix})
\end{multline*}
for all $X \in \M_{n\,k}(\F)$. On the other hand,
\begin{equation*}
T(X) = T(\begin{pmatrix}X[1|) - X[n|)\\ \vdots \\X[n-1|) - X[n|)\\ 0\end{pmatrix}) + T(\begin{pmatrix}X[n|)\\ \vdots \\ X[n|)\\ X[n|)\end{pmatrix})\\ = T(\begin{pmatrix}X[1|) - X[n|)\\ \vdots \\X[n-1|) - X[n|)\\ 0\end{pmatrix})
\end{equation*}
since $\begin{pmatrix}X[n|)\\ \vdots \\ X[n|)\\ X[n|)\end{pmatrix} \in W_{n\,k}$ and consequently $T(\begin{pmatrix}X[n|)\\ \vdots \\ X[n|)\\ X[n|)\end{pmatrix}) = 0$ by~\ref{lem:FromNKOddToNKEvenIfKerW:part4}.

\paragraph{\ref{lem:FromNKOddToNKEvenIfKerW:part2}} This part is a direct consequence from~\ref{lem:FromNKOddToNKEvenIfKerW:part5} and~\ref{lem:FromNKOddToNKEvenIfKerW:part6}. Indeed, 
\begin{multline*}
L^{+} \circ S \circ L^{-} \overset{\eqref{lem:FromNKOddToNKEvenIfKerW:eq}}{=\joinrel=} L^{+} \circ \left(L^{-} \circ T \circ L^{+} \right)\circ L^{-}\\
 = \left(L^{+} \circ L^{-} \circ T\right) \circ L^{+} \circ L^{-} \overset{\ref{lem:FromNKOddToNKEvenIfKerW:part5}}{=} T \circ L^{+} \circ L^{-} \overset{\ref{lem:FromNKOddToNKEvenIfKerW:part6}}{=} T.
\end{multline*}
\end{proof}

\begin{lemma}\label{lem:MatRepLift}Assume that  $|\F| > k$, $n > k$ and $n + k$ is odd. Let $S$ be a linear map on $\M_{(n-1)\,k}(\F)$ and $L$ be a linear map on $\M_{n\,k}(\F)$ defined by 
\[
T = L^{+}\circ S \circ L^{-}.
\] If there exist $A' \in \M_{(n-1)\,(n-1)}(\F)$ and $B' \in \M_{k\,k}(\F)$ satisfying the condition~\eqref{eq:TwoSideMul} such that 
\[
S(Y) = A'YB'\;\;\mbox{for all}\;\; Y \in \M_{(n-1)\,k}(\F),
\]
then there exist $A \in \M_{n\, n}(\F)$ and $B \in \M_{k\, k}(\F)$ satisfying the condition~\eqref{eq:TwoSideMul} such that
\[
T(X) = AXB\;\;\mbox{for all}\;\; X \in \M_{n\,k}(\F).
\]
\end{lemma}
\begin{proof}Since $A'$ and $B'$ satisfy  the condition~\eqref{eq:TwoSideMul}, then $S$ is a linear map preserving $\det_{(n-1)\,k}$ by Lemma~\ref{lem:TwoSidedMulPreservesDet}. Hence $T$ is a linear map preserving $\det_{n\,k}$ by Lemma~\ref{lem:LiftUpDetPres}.

The equality $T = L^{+} \circ S \circ L^{-}$ implies that
\[
T(X) = M^{+}(S(M^{-}X))
\]
for all $X \in \M_{n\,k}(\F)$. Recall there that $M^{+}$ and $M^{-}$ are defined in Lemma~\ref{claim:LminusMatrixForm} and Lemma~\ref{claim:LplusMatrixForm}. Therefore, \[
T(X) = M^{+}(A'(M^{-}X)B') = (M^{+}A'M^{-})XB' = AXB
\]
for all $X \in \M_{n\,k}(\F)$, where $A\in\M_{n\,n}(\F)$ and $B\in \M_{k\,k}(\F)$ are defined by
\[
A = M^{+}A'M^{-}, \quad B = B'.
\]
Since $T$ is a linear map preserving $\det_{n\,k}$ as we have shown above, then $A$ and $B$ satisfy the condition~\eqref{eq:TwoSideMul} by Lemma~\ref{lem:TwoSidedMulPreservesDet}. 
\end{proof}

\begin{theorem}\label{thm:MainTheoremCullisNKOdd}Assume that $|\F| > k \ge 4$, $n \ge k + 2$ and $n + k$ is odd. Let $T\colon \M_{n\,k} (\F) \to \M_{n\,k} (\F)$ be a linear map. Then $\det_{n\, k} (T(X)) = \det_{n\,k}(X)$ for all $X \in \M_{n\,k} (\F)$ if and only if  there exist $A \in \M_{n\, n}(\F)$ and $B \in \M_{k\, k}(\F)$ such that
\begin{equation}\label{thm:MainTheoremCullisNKOdd:eqq}
\det_{n\,k} \Bigl(A(|i_1,\ldots, i_k]\Bigr) \det_k \Bigl(B\Bigr) = (-1)^{i_1 + \ldots + i_k - 1 - \ldots - k}
\end{equation}
for all increasing sequences $1 \le i_1 < \ldots < i_k \le n $ and a linear map\linebreak $\phi\colon \M_{n\,k}(\F) \to W_{n\,k}$ such that
\begin{equation}\label{thm:MainTheoremCullisNKOdd:eq}
T(X) = AXB + \phi(X)
\end{equation}
for all $X \in \M_{n\,k} (\F).$

Here $W_{n\,k} \subseteq \M_{n\,k}(\F)$ is the space of matrices, all rows of which are equal, which is defined in Definition~\ref{def:WNK} and  is a radical of $\det_{n\,k}$ by Lemma~\ref{lem:NKOddRadDetNK}.
\end{theorem}
\begin{proof}
Let us show the sufficiency. Let $T$ be a linear map which satisfies the condition~\eqref{thm:MainTheoremCullisNKOdd:eq} for certain $A \in \M_{n\, n}(\F)$, $B \in \M_{k\, k}(\F)$ satisfying the condition~\eqref{thm:MainTheoremCullisNKOdd:eqq} and a linear map $\phi\colon \M_{n\,k}(\F) \to W_{n\,k}$. 

Lemma~\ref{lem:TwoSidedMulPreservesDet} implies that a linear map $T - \phi$ preserves $\det_{n\,k}$ . Therefore, $T = \left(T - \phi\right) + \phi$ preserves $\det_{n\,k}$ by Corollary~\ref{cor:AdditionLinearMapNPLusKOdd}. 

Let us prove the necessity. Assume that $T\colon \M_{n\,k} (\F) \to \M_{n\,k} (\F)$ is a linear map preserving $\det_{n\, k}$. Let $\phi\colon \M_{n\,k}(\F) \to W_{n\,k}$ be a linear map defined by
\[
\phi(X) = \begin{pmatrix}T[n|)\\\vdots\\ T[n|)\end{pmatrix},
\]
$\tilde{T}\colon \M_{n\,k} (\F) \to \M_{n\,k} (\F)$ be a linear map defined by $\tilde{T} = T - \phi$ and $S$ be a linear map on $\M_{(n-1)\,k}(\F)$  defined by $S = L^{-}\circ \tilde{T} \circ L^{+}$. 

The definition of $\tilde{T}$ implies that $\Img(\tilde{T}) \subseteq \M^{0}_{n\,k}(\F)$. In addition, $\tilde{T}$ preserves $\det_{n\,k}$ by Corollary~\ref{cor:AdditionLinearMapNPLusKOdd} applied to $T$ and $(-\phi)$. Thus, $\tilde{T}$ satisfies the conditions of Lemma~\ref{lem:FromNKOddToNKEvenIfKerW}. In particular, 
\begin{equation}\label{thm:MainTheoremCullisNKOdd:eq1}
\tilde{T} = L^{+}\circ S \circ L^{-}
\end{equation}
 by Lemma~\ref{lem:FromNKOddToNKEvenIfKerW}\ref{lem:FromNKOddToNKEvenIfKerW:part2}.

Since $\tilde{T}$ preserves $\det_{n\,k}$, then $S$ preserves $\det_{(n-1)\,k}$ by Lemma~\ref{lem:LiftDownDetPres}.\linebreak Lemma~\ref{lem:MainTheoremEven} implies that there exist $A' \in \M_{(n-1)\,(n-1)}(\F)$ and $B' \in \M_{k\,k}(\F)$ satisfying the condition~\eqref{eq:TwoSideMul} such that 
\begin{equation}\label{thm:MainTheoremCullisNKOdd:eq2}
S(Y) = A'YB'\;\;\mbox{for all}\;\; Y \in \M_{(n-1)\,k}(\F).
\end{equation}

Thus, the equalities~\eqref{thm:MainTheoremCullisNKOdd:eq1} and~\eqref{thm:MainTheoremCullisNKOdd:eq2} allow us to conclude from Lemma~\ref{lem:MatRepLift} that there exist $A \in \M_{n\, n}(\F)$ and $B \in \M_{k\, k}(\F)$ satisfying the condition~\eqref{eq:TwoSideMul} such that 
$$\tilde{T}(X) = AXB\;\;\mbox{for all}\;\; X \in \M_{n\,k}(\F).$$
Therefore,
\[
T(X) = \tilde{T}(X) + \phi(X) = AXB + \phi(X)
\]
for all $X \in \M_{n\,k} (\F).$
\end{proof}

\begin{remark}If $T$ is a linear map on $\M_{n\,k}$ preserving $\det_{n\,k}$, then the representation~\eqref{thm:MainTheoremCullisNKOdd:eq} may be not unique. For example,
\[
I_nXI_k + \begin{pmatrix}x_{1\,1} & \cdots & x_{1\,k} \\ \vdots & \ddots & \vdots\\ x_{1\,1} & \cdots & x_{1\,k}\end{pmatrix} = AXI_k\;\;\mbox{for all}\;\; X \in \M_{n\,k}(\F),
\]
where $A  = E_{1\,1} + \ldots + E_{n\,n} + E_{1\,n} + \ldots + E_{n\,n}$, and the pair of matrices $(A, I_k)$ satisfies the condition~\eqref{thm:MainTheoremCullisNKOdd:eqq}.
\end{remark}

\begin{remark}There exist linear maps of the form~\eqref{thm:MainTheoremCullisNKOdd:eq} which do not admit a representation $X \mapsto A'XB'$ for any $A' \in \M_{n\,n}$, $B' \in \M_{k\,k}$. For example a linear map $T'\colon \M_{n\,k}(\F) \to \M_{n\,k}(\F)$ defined by  
\[
T'(X) =  X + \begin{pmatrix}-x_{1\,1} & 0 & \cdots & 0\\ \vdots & \vdots & \ddots & \vdots\\ -x_{1\,1} & 0 & \cdots & 0
\end{pmatrix}\;\;\mbox{for all}\;\; X \in \M_{n\,k}(\F)
\]
has the form~\eqref{thm:MainTheoremCullisNKOdd:eq} but increases a matrix rank because
\[
\rk\left(E_{1\,1} + E_{1\,2}\right) = 1
\]
and
\[
\rk\left(T'(E_{1\,1} + E_{1\,2})\right) = \rk\left(-E_{2\,1} + \ldots + -E_{n\,1} + E_{1\,2}\right) = 2.
\]
This implies that there do not exist $A' \in \M_{n\,n}$, $B' \in \M_{k\,k}$ such that 
\begin{equation}\label{remeq1}
T'(X) = A'X'B\;\;\mbox{for all}\;\; X \in \M_{n\,k}(\F)
\end{equation}
since every map of the form~\eqref{remeq1} cannot increase the rank of an argument.
\end{remark}



\section*{Declaration of competing interest}

The authors declare that they have no known competing financial interests or personal relationships that could have appeared to influence the work reported in this paper.

\section*{Acknowledgements}

The research of the second author was supported by the scholarship of the Center for Absorption in Science, the Ministry for Absorption of Aliyah, the State of Israel.

\section*{Data availability}

No data was used for the research described in the article.

\bibliographystyle{plain}
\bibliography{cullisgeneral}




\end{document}
